\section{Loop recognition, lawful ledgers, and curvature}

\subsection{Local ledger theorem}

Let $D$ be an oriented smooth domain, let $V_0$ be a finite-dimensional vector space, and let $\omega,\ell\in\Omega^1(D;V_0)$.  The form $\omega$ is the measured response and $\ell$ is a prior lawful local ledger.  Define the open loop residue
\begin{equation}
G_{\omega,\ell}(\gamma)
:=\oint_\gamma(\omega-\ell).
\label{eq:open-loop-residue}
\end{equation}

\begin{definition}[Loop Recognition Kernel]
The loop Recognition Kernel is
\[
\cK_{\omega,\ell}
:=\{\gamma\text{ closed}:G_{\omega,\ell}(\gamma)=0\}.
\]
\end{definition}

\begin{theorem}[Recognition--Stokes theorem]
If $A\subset D$ is an oriented piecewise smooth two-chain, then
\begin{equation}
G_{\omega,\ell}(\partial A)
=\int_A\dd(\omega-\ell).
\label{eq:recognition-stokes}
\end{equation}
If $D$ is simply connected, the following are equivalent:
\begin{enumerate}[label=(\roman*)]
\item every closed loop belongs to $\cK_{\omega,\ell}$;
\item $\dd(\omega-\ell)=0$;
\item there exists a $V_0$-valued potential $\Phi$ with $\omega-\ell=\dd\Phi$.
\end{enumerate}
\end{theorem}

\begin{proof}
Equation \eqref{eq:recognition-stokes} is Stokes' theorem.  The implication (iii)$\Rightarrow$(i) is the fundamental theorem for line integrals, and (i)$\Rightarrow$(ii) follows by applying \eqref{eq:recognition-stokes} to arbitrarily small coordinate rectangles.  For (ii)$\Rightarrow$(iii), fix a base point and define $\Phi(x)$ by integrating $\omega-\ell$ along a path from the base point to $x$.  Closedness and simple connectivity make the integral path independent, and differentiation gives $\dd\Phi=\omega-\ell$.
\end{proof}

\begin{remark}[History-dependent memory]
A genuinely history-dependent ledger need not be representable by a one-form on the equilibrium state manifold.  In that case the state must be enlarged by internal or memory variables.  Stokes' theorem is then applied on the enlarged state space.  Subtracting a path-dependent term after observing the loop, without declaring it as prior state data, is not a lawful closure argument.
\end{remark}

\subsection{Return classes}

The residue accounting leads to four mutually distinguishable classes:
\begin{enumerate}[label=\textbf{Class \Roman*:}]
\item exact raw return, $\oint\omega=0$;
\item lawfully transported return, $\oint\omega=\oint\ell$;
\item memory-bearing return on an enlarged state space, with the memory component nonzero but fully ledgered;
\item open obstruction, where $G_{\omega,\ell}$ exceeds the declared uncertainty or tolerance.
\end{enumerate}

These are classifications, not aliases.  A nonzero geometric loop integral is not automatically irreversible.

\subsection{Onsager covariance and calibration}

Let $X,J\in\R^m$ be thermodynamic forces and conjugate fluxes, with entropy-production rate in the linear irreversible-thermodynamic convention \cite{Onsager1,Onsager2,Casimir}
\[
\sigma=J^{\mathsf T}X,
\qquad
J=LX.
\]

\begin{theorem}[Covariance of reciprocal linear response]
Let $B$ be invertible and transform the dual variables by
\[
J'=BJ,
\qquad
X'=B^{-\mathsf T}X.
\]
Then
\[
\sigma=J'^{\mathsf T}X',
\qquad
L'=BLB^{\mathsf T}.
\]
If $L=L^{\mathsf T}\ge0$, then $L'=L'^{\mathsf T}\ge0$.
\end{theorem}

\begin{proof}
The pairing is invariant because
\[
J'^{\mathsf T}X'=J^{\mathsf T}B^{\mathsf T}B^{-\mathsf T}X=J^{\mathsf T}X.
\]
Since $X=B^{\mathsf T}X'$, one has
\[
J'=BLB^{\mathsf T}X'.
\]
Symmetry and positive semidefiniteness are preserved by congruence.
\end{proof}

\begin{corollary}[Frame asymmetry is not automatically nonreciprocity]
A coefficient matrix that becomes asymmetric after a non-dual rescaling of forces or fluxes does not constitute an invariant violation of Onsager reciprocity.  A physical nonreciprocity claim must be made in a conjugate flux--force frame and must include the relevant Onsager--Casimir parity assumptions.
\end{corollary}

\begin{principle}[Curvature before entropy production]
The equality
\[
\int_A\dd(\omega-\ell)=\int_A q_\sigma\,\dd A
\]
is not a geometric identity unless the constitutive density $q_\sigma$ and its units, sign, time orientation, and calibration are independently supplied.  Geometry determines the left-hand residue; thermodynamics determines whether and how it represents dissipation.
\end{principle}

\section{Finite Cut--Flow--Jet Recognition theorem}
\label{sec:cut-flow-capstone}

\subsection{Finite graded carrier and jet raiser}

Let
\[
\cJ_N=\bigoplus_{n=0}^N H_n
\]
be a finite-dimensional graded Hilbert space, with inclusions $\iota_n:H_n\to\cJ_N$ and projections $\pi_n$.  A raising operator $\mathbb R$ satisfies
\[
\mathbb R(H_n)\subseteq H_{n+1},
\qquad
\mathbb R(H_N)=0.
\]
Hence $\mathbb R^{N+1}=0$.

For a response bundle with connection $\nabla$, the holonomic jet realization is obtained when
\begin{equation}
\mathbb R^n\iota_0\Lambda
=\iota_n\Sym\nabla^n\Lambda,
\qquad 0\le n\le N.
\label{eq:holonomic-raiser}
\end{equation}
This condition is part of the datum; a level shift by the identity alone certifies only combinatorial tower coefficients, not covariant derivatives.

Let $\mathbb T$ preserve the grading and set
\[
\mathbb G=\mathbb T+\mathbb R.
\]

\begin{theorem}[Finite flow-generated jet tower]
If $[\mathbb T,\mathbb R]=0$, then
\begin{equation}
\e^{t\mathbb G}=\e^{t\mathbb T}\e^{t\mathbb R}
\label{eq:flow-factorization}
\end{equation}
and, for a holonomic seed,
\begin{equation}
\pi_n\e^{t\mathbb G}\iota_0\Lambda
=\frac{t^n}{n!}\,
\pi_n\e^{t\mathbb T}\iota_n\Sym\nabla^n\Lambda.
\label{eq:jet-flow-formula}
\end{equation}
If the commutator is nonzero, then the first factorization defect is
\begin{equation}
\e^{t\mathbb T}\e^{t\mathbb R}
=\exp\left(
 t(\mathbb T+\mathbb R)
 +\frac{t^2}{2}[\mathbb T,\mathbb R]
 +O(t^3)
\right).
\label{eq:jet-bch}
\end{equation}
\end{theorem}

\begin{proof}
Because $\mathbb R^{N+1}=0$,
\[
\e^{t\mathbb R}=\sum_{k=0}^N\frac{t^k}{k!}\mathbb R^k.
\]
Commutation gives \eqref{eq:flow-factorization}; projection to layer $n$ and \eqref{eq:holonomic-raiser} give \eqref{eq:jet-flow-formula}.  Equation \eqref{eq:jet-bch} is the Baker--Campbell--Hausdorff expansion, valid near $t=0$ for finite matrices.
\end{proof}

\subsection{Self-adjoint cut and unique grading}

Let $\mathsf J$ be a self-adjoint unitary involution,
\begin{equation}
\mathsf J^*=\mathsf J=\mathsf J^{-1}.
\label{eq:selfadjoint-cut}
\end{equation}

\begin{theorem}[Unique cut grading]
Every operator $A$ on $\cJ_N$ decomposes uniquely as
\[
A=A_{\mathrm e}+A_{\mathrm o},
\qquad
A_{\mathrm e}=\frac12(A+\mathsf J A\mathsf J),
\qquad
A_{\mathrm o}=\frac12(A-\mathsf J A\mathsf J),
\]
with
\[
\mathsf J A_{\mathrm e}\mathsf J=A_{\mathrm e},
\qquad
\mathsf J A_{\mathrm o}\mathsf J=-A_{\mathrm o}.
\]
\end{theorem}

\begin{proof}
Direct substitution proves the parity relations and $A=A_{\mathrm e}+A_{\mathrm o}$.  If $A=B_{\mathrm e}+B_{\mathrm o}$ is another such decomposition, applying the two parity projectors gives $B_{\mathrm e}=A_{\mathrm e}$ and $B_{\mathrm o}=A_{\mathrm o}$.
\end{proof}

\subsection{Cut loop and commutator curvature}

Set $U_t=\e^{t\mathbb G}$ and define
\[
\mathscr H_{\mathsf J}(t)
:=\mathsf J U_t\mathsf J U_t.
\]

\begin{theorem}[Cut-loop expansion]
For sufficiently small $t$, the principal logarithm is defined and
\begin{equation}
\log\mathscr H_{\mathsf J}(t)
=2t\mathbb G_{\mathrm e}
+t^2[\mathbb G_{\mathrm e},\mathbb G_{\mathrm o}]
+O(t^3).
\label{eq:cut-loop-expansion}
\end{equation}
Thus the cut-even generator is the first seam-preserving term, while
\[
\cR_{\mathrm{seam}}
:=[\mathbb G_{\mathrm e},\mathbb G_{\mathrm o}]
\]
is the first noncommutative cut-loop curvature coefficient.
\end{theorem}

\begin{proof}
Since
\[
\mathsf J\mathbb G\mathsf J
=\mathbb G_{\mathrm e}-\mathbb G_{\mathrm o},
\]
one has
\[
\mathscr H_{\mathsf J}(t)
=\e^{t(\mathbb G_{\mathrm e}-\mathbb G_{\mathrm o})}
 \e^{t(\mathbb G_{\mathrm e}+\mathbb G_{\mathrm o})}.
\]
The BCH formula, with the remainder taken in operator norm, gives the sum $2t\mathbb G_{\mathrm e}$ and one half of the commutator
\[
\frac12 t^2
[\mathbb G_{\mathrm e}-\mathbb G_{\mathrm o},
 \mathbb G_{\mathrm e}+\mathbb G_{\mathrm o}]
=t^2[\mathbb G_{\mathrm e},\mathbb G_{\mathrm o}],
\]
which proves \eqref{eq:cut-loop-expansion}.
\end{proof}

\begin{corollary}[Bilateral reconstruction for a cut-odd generator]
If $\mathsf J\mathbb G\mathsf J=-\mathbb G$, then
\[
\mathsf J U_t\mathsf J=U_{-t}.
\]
With
\[
E_t=\frac12(U_t+U_{-t}),
\qquad
O_t=\frac12(U_t-U_{-t}),
\]
one has
\[
U_t=E_t+O_t,
\qquad
U_{-t}=E_t-O_t,
\qquad
E_t^2-O_t^2=\id.
\]
\end{corollary}

\begin{proof}
Conjugation of the exponential gives the first identity.  The reconstruction formulas are definitions, and
\[
E_t^2-O_t^2
=\frac14\bigl[(U_t+U_{-t})^2-(U_t-U_{-t})^2\bigr]
=\frac12(U_tU_{-t}+U_{-t}U_t)=\id.
\]
\end{proof}

\subsection{Cut and adjoint bi-grading}

\begin{theorem}[Cut--adjoint bi-grading]
For $\sigma,\tau\in\{+1,-1\}$, define
\begin{equation}
A^{\sigma,\tau}
:=\frac14\left(
A+\sigma\mathsf J A\mathsf J
+\tau A^*
+\sigma\tau\mathsf J A^*\mathsf J
\right).
\label{eq:bigrading}
\end{equation}
Then
\[
A=\sum_{\sigma,\tau}A^{\sigma,\tau},
\qquad
\mathsf J A^{\sigma,\tau}\mathsf J=\sigma A^{\sigma,\tau},
\qquad
(A^{\sigma,\tau})^*=\tau A^{\sigma,\tau}.
\]
The four sectors are pairwise orthogonal in the real Hilbert--Schmidt inner product
\[
(A,B)_{\mathrm{HS},\R}:=\operatorname{Re}\tr(A^*B).
\]
\end{theorem}

\begin{proof}
Viewed on the underlying real vector space of operators with inner product $(\cdot,\cdot)_{\mathrm{HS},\R}$, cut conjugation and adjoint are commuting self-adjoint involutions.  Formula \eqref{eq:bigrading} is the product of their real spectral projectors.  The stated parity properties and real orthogonality follow.  The real qualification is essential: self-adjoint and skew-adjoint operators need not be orthogonal for the complex-valued Hilbert--Schmidt pairing.
\end{proof}

\begin{remark}[Required cut hypothesis]
The self-adjointness in \eqref{eq:selfadjoint-cut} is essential.  An arbitrary algebraic involution $\mathsf J^2=\id$ need not commute with the adjoint operation, and then \eqref{eq:bigrading} need not have the claimed adjoint parity.
\end{remark}

\subsection{Closure and the clock-free Eye}

Let $K=K^*$ and $U_t=\e^{-itK}$.  A single sampling time can contain stroboscopic fixed modes with $\e^{-it_0\lambda}=1$ even when $\lambda\ne0$.  The all-time average removes this ambiguity; this is the continuous mean-ergodic projection in spectral form \cite{ReedSimon}.

\begin{theorem}[Clock-free Eye]
The strong limit
\begin{equation}
P_{\Eye}
=\operatorname*{s-lim}_{T\to\infty}
\frac1T\int_0^T\e^{-itK}\,\dd t
=E_K(\{0\})
\label{eq:eye}
\end{equation}
exists and is the spectral projection onto $\Ker K$.
\end{theorem}

\begin{proof}
For $x\in\cH$, the spectral theorem gives
\[
\frac1T\int_0^T\e^{-itK}x\,\dd t
=\int_\R
\left(\frac1T\int_0^T\e^{-it\lambda}\,\dd t\right)
\dd E_K(\lambda)x.
\]
The scalar factor tends to $1$ at $\lambda=0$ and to $0$ otherwise, while its modulus is at most one.  Dominated convergence for the spectral measure gives the strong limit $E_K(\{0\})x$.
\end{proof}

\subsection{Representation invariance}

\begin{proposition}[Unitary invariance]
Let $W=\bigoplus_{n=0}^N W_n:\cJ_N\to\cJ_N'$ be a grade-preserving unitary.  Transport the inclusions, projections, seed, and operators by
\[
\iota_n'=W\iota_nW_n^*,\quad
\pi_n'=W_n\pi_nW^*,\quad
\Lambda'=W_0\Lambda,
\]
\[
\mathbb R'=W\mathbb RW^*,\quad
\mathbb T'=W\mathbb TW^*,\quad
\mathsf J'=W\mathsf JW^*,\quad
K'=WKW^*,
\]
and transport observations and scalar targets by
\[
\cA'=\cA W^*,\qquad
\Pi'=\Pi W^*,\qquad
L'=LW^*.
\]
Then cut grading, finite jet coefficients, cut-loop curvature, the Eye projection, blind-quotient dimension, and scalar decoder burden are preserved.
\end{proposition}

\begin{proof}
All operator identities are preserved by unitary conjugation.  Grade preservation gives
\[
\pi_n'\e^{t\mathbb G'}\iota_0'\Lambda'
=W_n\pi_n\e^{t\mathbb G}\iota_0\Lambda.
\]
Kernels are carried bijectively by $W$ and quotient dimensions are unchanged.  Spectral projections obey
\[
E_{K'}(\{0\})=WE_K(\{0\})W^*.
\]
The numerator and denominator in \eqref{eq:decoder-burden} are unchanged after the substitution $x'=Wx$.
\end{proof}
